% page 261 du fac-simile (image p-260 du scan)
% bloc 182.9 x 242.0 mm ; marge gauche 19.1 mm
\pgc{-2.540mm}{0.000mm}{
\l{\cel{59}253}
\l{}
\l{\sou{kajo}. I. Quaza chambro, garnisita ye stangi ek ligno, ek fero, en qua}
\l{\cel{3}onu enklozas animali ne-amansita por gardar li en menajerio o por}
\l{\cel{3}transportar li.\cel{2}II. Quaza chambreto portebla, garnisita ye latun-}
\l{\cel{3}fili, en qua onu retenas enklozite uceli. - EFI.}
\l{}
\l{\sou{kakar}. (\sou{netrans.})\cel{2}Vakuigar sua ventro de la grosa exkrementi. - DEFIS.}
\l{}
\l{\sou{kakao}. (\sou{bot.}) Grano kontenata da la frukto di arbusto ek la familio}
\l{\cel{3}\textquotedbl{}malvacei\textquotedbl{} (L.\cel{1}\sou{theobroma cacao}) e qua, grilite e triturite kun}
\l{\cel{3}sukro, donas chokolado. - DEFIRS.}
\l{}
\l{\sou{kakatuo.}\cel{2}(\sou{zool.}) Speco di papagayo penachoza, ucelo klimera, di qua}
\l{\cel{3}la plumaro esas bele blanka. - L.\cel{1}\sou{piictolophus}. - DEFIRS.}
\l{}
\l{\sou{kakexio.}\cel{2}(\sou{patol.}) Stando maladatra, quan karakterizas deperiso ge-}
\l{\cel{3}nerala. - DEFS.}
\l{}
\l{\sou{kakofonio}.\cel{2}I. Konsonanco qua lezas la oreli\cel{2}- II. (\sou{muziko)}\cel{1}.}
\l{\cel{3}Asemblajo ek son-uni ne-akordanta. - DEFIS.}
\l{}
\l{\sou{kaktuso}. (\sou{bot.)}\cel{2}Planto di la familio \textquotedbl{}kaktei\textquotedbl{}, qua donas frukto}
\l{\cel{3}laxigiva, figatra. - L.\cel{1}\sou{cactus}. - DEFIRS.}
\l{}
\l{\sou{kalo}. (\sou{patol}.)Dikesko od hardesko di la pelo, produktita da fricion-}
\l{\cel{3}ado, ye la palmo, ye la plando, sur la pedo-fingri. - EFIS.}
\l{}
\l{\sou{kalamo}. Fragmento de kano, quan la antiqui uzis por skribar.-DFIS.}
\l{}
\l{\sou{kalamino}. (\sou{kemio}). Nomo vulgara di la hidro-silikato di zinko.}
\l{\cel{3}\sou{Simbolo kemiala}\cel{1}: Si O5 Z2 H2. - DEFIS.}
\l{}
\l{\sou{kalamito}. Rezino qualeso-inferiora,quan onu rekoltas de la kano-}
\l{\cel{3}stipi. - L.\cel{1}\sou{styrax calamit}a. - DEFIS.}
\l{}
\l{\sou{kalamitato.}\cel{3}Granda desfortuno qua unafoye frapas tre multa}
\l{\cel{3}personi e mem populo. - DEFIS.}
\l{}
\l{\sou{kalandrar}.\cel{2}(\sou{trans.}) Pasigar sub cilindro glata la stofi,drapi,}
\l{\cel{3}teli, e c. quin onu volas briligar, muarizar, e c. - DEFS.}
\l{}
\l{\sou{kalcedono}.\cel{2}(\sou{mineral.)}\cel{2}Varietato di agato, di qua la diafaneso}
\l{\cel{3}esas laktatra. - DEFIRS.}
\l{}
\l{\sou{kalcio}. (\sou{kemi}o).\cel{2}Korpo simpla, metala, flave blanka. -\cel{1}\sou{Simbolo kemiala}:}
\l{\cel{3}Ca.\cel{2}- DEFIS.}
\l{}
\l{\sou{kalcinar}. (\sou{trans.})\cel{2}Transformar metalo aden oxido, per submisar olu}
\l{\cel{3}a la efiko da fairo intensa, en aero libera. - DEFIS.}
\l{}
\l{\sou{kaldiero.}\cel{2}Vazo klozita (apartenanta, generale, a vapormashino) en qua}
\l{\cel{3}onu boliigas aquo por obtenar vaporo. - EFIS.}
\l{}
\l{\sou{kaldrono.}\cel{2}Mikra recipiento, ordinare de kupro, en qua onu koquas}
\l{\cel{3}la nutrivi o boliigas aquo. - EFIS.}
}
